\documentclass[a4paper,11pt]{article}

% Packages
\usepackage[utf8]{inputenc}
\usepackage[T1]{fontenc}
\usepackage[french]{babel}
\usepackage{amsmath, amssymb, amsthm}
\usepackage{geometry}
\usepackage{xcolor}
\usepackage{tcolorbox}
\usepackage{enumitem}
\usepackage{fancyhdr}
\usepackage{array}
\usepackage{booktabs}
\usepackage{listings}
\usepackage{hyperref}

% Configuration de la page
\geometry{margin=2.5cm}
\pagestyle{fancy}
\fancyhf{}
\fancyhead[L]{Mathématiques Discrètes - TDs}
\fancyhead[R]{Licence 1 Informatique}
\fancyfoot[C]{\thepage}

% Environnements personnalisés
\tcbuselibrary{theorems,skins,breakable}

\newtcolorbox{exercisebox}[1]{
    colback=blue!5!white,
    colframe=blue!75!black,
    fonttitle=\bfseries,
    title=#1,
    breakable
}

\newtcolorbox{solutionbox}{
    colback=green!5!white,
    colframe=green!75!black,
    fonttitle=\bfseries,
    title=Solution,
    breakable
}

\newtcolorbox{notebox}{
    colback=orange!5!white,
    colframe=orange!75!black,
    fonttitle=\bfseries,
    title=Note importante,
    breakable
}

\newtcolorbox{remarkbox}{
    colback=purple!5!white,
    colframe=purple!75!black,
    fonttitle=\bfseries,
    title=Remarque,
    breakable
}

% Configuration listings pour pseudo-code
\lstset{
    basicstyle=\ttfamily\small,
    keywordstyle=\bfseries,
    commentstyle=\itshape\color{gray},
    frame=single,
    breaklines=true,
    showstringspaces=false
}

% Titre
\title{\textbf{Mathématiques Discrètes} \\ 
\large Travaux Dirigés - Licence 1 Informatique}
\author{CERI - Avignon Université}
\date{Hiver 2026}

\begin{document}

\maketitle
\thispagestyle{empty}

\newpage
\tableofcontents
\newpage

% ============================================================================
% TD1 - MANIPULER LA LOGIQUE
% ============================================================================

\part{TD1 - Manipuler la logique}

\section*{Objectif}
L'objectif de ce TD est de manipuler les concepts de base de la logique afin de les prendre en main pour les maîtriser.

\section{Formules et priorité}

\begin{exercisebox}{Exercice 1}
Ajoutez des parenthèses aux formules à priorités ci-dessous afin de les transformer en formules strictes équivalentes.

\begin{enumerate}
    \item $p \vee q \vee r \vee s \vee t$
    \item $p \wedge q \wedge r \wedge s \wedge t$
    \item $p \Leftrightarrow \neg q \vee r$
    \item $p \vee q \Rightarrow r \wedge s$
    \item $p \vee q \Rightarrow r \Leftrightarrow s$
    \item $p \vee q \wedge r \Rightarrow \neg s$
    \item $p \Rightarrow r \wedge s \Rightarrow t$
    \item $p \vee q \wedge s \vee t$
    \item $p \wedge q \Leftrightarrow \neg r \vee s$
    \item $\neg p \wedge q \vee r \Rightarrow s \Leftrightarrow t$
\end{enumerate}
\end{exercisebox}

\section{Logique et français}

\begin{exercisebox}{Exercice 2}
Traduisez les phrases en français ci-dessous en formules logiques en précisant quelles variables sont utilisées pour quelles propositions.

\begin{enumerate}
    \item Il est faux de dire que nous sommes des géants.
    \item Nous sommes inscrits à l'université mais pas à AMU.
    \item Les étudiants qui ont suivi des cours de mathématiques ou d'informatique peuvent s'inscrire.
    \item Je vais au ski cet hiver ou je voyage à l'étranger cet été, mais pas les deux.
    \item Il est nécessaire et suffisant d'être majeur pour avoir le droit de voter.
    \item Je suis riche seulement si je gagne à la loterie.
    \item Si le temps est clément et que je suis en forme, alors je vais me balader.
    \item Aujourd'hui n'est pas un jour férié.
    \item Je suis des cours de mathématiques et des cours d'informatique.
    \item S'il pleut, j'irai faire les courses. S'il ne pleut pas, j'irai faire les courses quand même, et j'irai en plus me balader à la plage.
\end{enumerate}
\end{exercisebox}

\section{Implication logique et français}

\begin{exercisebox}{Exercice 3}
Pour chacune des phrases en français suivantes, donnez la phrase correspondant à la proposition inverse, la proposition réciproque et la proposition contraposée. Précisez quelles phrases ont un sens équivalent.

\begin{enumerate}
    \item S'il pleut, je vais faire les courses.
    \item J'ai mon diplôme seulement si je valide mon année.
    \item Si le temps est clément et que je suis en forme, alors je vais me balader.
    \item Si le temps est clément ou que je suis en forme, alors je vais me balader.
    \item Si le temps est clément, alors je vais me balader et manger une glace.
    \item Si le temps est clément, alors je vais me balader ou manger une glace.
    \item Si le temps est clément et que je suis en forme, alors je vais me balader et manger une glace.
    \item Si le temps est clément et que je suis en forme, alors je vais me balader ou manger une glace.
    \item Si le temps est clément ou que je suis en forme, alors je vais me balader et manger une glace.
    \item Si le temps est clément ou que je suis en forme, alors je vais me balader ou manger une glace.
\end{enumerate}
\end{exercisebox}

\section{Manipulation de formules}

\begin{exercisebox}{Exercice 4}
Montrez les équivalences suivantes en utilisant les équivalences remarquables :

\begin{enumerate}
    \item $\neg(a \Rightarrow b) \equiv a \wedge \neg b$
    \item $a \wedge b \Rightarrow a \equiv \top$
    \item $a \Rightarrow a \vee b \equiv \top$
    \item $x \Rightarrow (y \wedge z) \equiv (x \Rightarrow y) \wedge (x \Rightarrow z)$
    \item $x \Rightarrow (y \vee z) \equiv (x \Rightarrow y) \vee (x \Rightarrow z)$
    \item $(a \vee b) \Rightarrow (c \wedge d) \equiv ((a \Rightarrow c) \wedge (a \Rightarrow d)) \wedge ((b \Rightarrow c) \wedge (b \Rightarrow d))$
    \item $(a \vee b) \Rightarrow (x \wedge y) \equiv (a \Rightarrow (x \wedge y)) \wedge (b \Rightarrow (x \wedge y))$
    \item $(a \wedge b) \Rightarrow (x \vee y) \equiv (a \Rightarrow x) \vee (b \Rightarrow y)$
    \item $(x \Rightarrow y) \Rightarrow x \equiv x \vee y$
    \item $\neg(x \Rightarrow y) \equiv x \Leftrightarrow \neg y$
\end{enumerate}
\end{exercisebox}

\section{Procédure booléenne}

\begin{exercisebox}{Exercice 5}
On fournit le pseudo-code de la procédure \texttt{booleenne1}.

\begin{lstlisting}
Procedure booleenne1()
    x <- Entier aleatoire entre 1 et 10
    y <- Entier aleatoire entre 1 et 10
    Si x > y Alors
        Renvoyer FAUX
    Sinon
        Renvoyer VRAI
    FinSi
FinProcedure
\end{lstlisting}

\begin{enumerate}
    \item Dans quel cas cette procédure renvoie-t-elle la valeur FAUX ?
    \item Dans quel cas cette procédure renvoie-t-elle la valeur VRAI ?
    \item Proposez une procédure \texttt{booleenne2} équivalente.
\end{enumerate}
\end{exercisebox}

\section{Quelques boucles}

\begin{exercisebox}{Exercice 6}
\textbf{Question 1 :} On fournit le pseudo-code de la procédure \texttt{tirage1}. Que peut-on dire une fois cette exécution terminée ? Utilisez une formalisation logique afin de répondre.

\begin{lstlisting}
Procedure tirage1()
    a <- Valeur du tirage d'une piece a pile ou face
    b <- Valeur du tirage d'une piece a pile ou face
    TantQue (a.estPile()&b.estPile())||
            (a.estFace()&b.estFace()) Faire
        a <- Valeur du tirage d'une piece a pile ou face
        b <- Valeur du tirage d'une piece a pile ou face
    FinTantQue
FinProcedure
\end{lstlisting}

\textbf{Question 2 :} Même question pour \texttt{tirage2}.

\begin{lstlisting}
Procedure tirage2()
    a <- Valeur du tirage d'une piece a pile ou face
    b <- Valeur du tirage d'une piece a pile ou face
    TantQue (!a.estPile()&b.estPile())||
            (a.estFace()&!b.estFace()) Faire
        a <- Valeur du tirage d'une piece a pile ou face
        b <- Valeur du tirage d'une piece a pile ou face
    FinTantQue
FinProcedure
\end{lstlisting}

\textbf{Question 3 :} Peut-on simplifier ces programmes ?
\end{exercisebox}

\section{Quelques conditionnelles}

\begin{exercisebox}{Exercice 7}
On fournit le pseudo-code de la procédure \texttt{conditionnelle1}. On donne les variables logiques et propositions suivantes :
\begin{itemize}
    \item $a$ : $x \leq y$, connaissant $x$ et $y$.
    \item $b$ : $x = y + 1$, connaissant $x$ et $y$.
    \item $t$ : On affiche 'Titi'
    \item $o$ : On affiche 'Toto'
    \item $u$ : On affiche 'Tutu'
\end{itemize}

\begin{lstlisting}
Procedure conditionnelle1()
    x <- Valeur du tirage d'un de a 6 faces
    y <- Valeur du tirage d'un de a 6 faces
    Si x > y Alors
        Afficher "Titi"
    SinonSi x = y + 1 Alors
        Afficher "Toto"
    Sinon
        Afficher "Tutu"
    FinSi
FinProcedure
\end{lstlisting}

\begin{enumerate}
    \item Exprimez en fonction de $a$ et $b$ le cas où on affiche 'Titi'.
    \item Exprimez en fonction de $a$ et $b$ le cas où on affiche 'Toto'.
    \item Exprimez en fonction de $a$ et $b$ le cas où on affiche 'Tutu'.
    \item Complétez les conditions dans la procédure suivante.
\end{enumerate}
\end{exercisebox}

\newpage

% ============================================================================
% TD2 - RÉSOLUTION
% ============================================================================

\part{TD2 - Résolution}

\section*{Objectif}
L'objectif de ce TD est de mettre en pratique les notions de résolution.

\section{Formes normales réduites}

\begin{exercisebox}{Exercice 8}
À partir des formules, donnez la forme normale conjonctive et réduire cette formule. Donnez les contre-modèles de la formule.

\begin{enumerate}
    \item $y \Leftrightarrow (x \wedge y)$
    \item $(x \Leftrightarrow (y \vee z)) \wedge (z \Rightarrow (x \vee y)) \wedge \neg(y \vee \neg x)$
    \item $(x \wedge y) \vee (\neg x \wedge y \wedge \neg z)$
    \item $(\neg p \Rightarrow q) \wedge ((p \Rightarrow r) \vee (q \Rightarrow p)) \wedge (r \Rightarrow p) \wedge (\neg q \Rightarrow \neg(p \wedge q)) \wedge (r \Rightarrow (p \vee q)) \wedge ((p \Rightarrow r) \vee (r \Rightarrow p)) \wedge (\neg q \Rightarrow \neg r)$
    \item $(x \vee y \vee z) \wedge (\neg x \Rightarrow y) \wedge (z \Leftrightarrow y)$
\end{enumerate}
\end{exercisebox}

\section{Identifier des inférences}

\begin{exercisebox}{Exercice 9}
Pour chacun des cas, utilisez un formalisme logique pour traduire les hypothèses et la conclusion et précisez la règle d'inférence utilisée pour obtenir cette conclusion à partir des hypothèses.

\begin{enumerate}
    \item Alice étudie les mathématiques. On en déduit qu'Alice étudie les mathématiques ou l'informatique.
    \item Bob étudie les mathématiques et l'informatique. On en déduit que Bob étudie les mathématiques.
    \item S'il pleut alors la piscine est fermée. Aujourd'hui il pleut. On en déduit que la piscine est fermée.
    \item S'il neige, l'université est fermée. L'université est ouverte aujourd'hui. On en déduit qu'il ne neige pas.
    \item Quand je me baigne, je reste trop longtemps au soleil. Si je reste au soleil trop longtemps, alors j'ai des coups de soleil. On en déduit que si je vais me baigner, alors j'aurai des coups de soleil.
\end{enumerate}
\end{exercisebox}

\section{Déductions}

\begin{exercisebox}{Exercice 10}
Pour chacun des cas, utilisez les règles d'inférence afin de détailler l'argumentation permettant d'obtenir la conclusion donnée.

\begin{enumerate}
    \item Déduire $c$ à partir des hypothèses : $\neg d$, $(\neg q \vee b)$, $(b \Rightarrow d)$, $(a \vee c)$
    
    \item Déduire $e$ à partir des hypothèses : $(\neg a \wedge b)$, $(d \Rightarrow c)$, $(\neg e \Rightarrow d)$, $(c \Rightarrow a)$
    
    \item Déduire $a \wedge b$ à partir des hypothèses : $(\neg e \Rightarrow \neg d)$, $(\neg f \Rightarrow d)$, $(a \vee c)$, $(\neg e \wedge \neg c)$
    
    \item Déduire $d$ à partir des hypothèses : $(a \Rightarrow b)$, $\neg c$, $(a \vee d)$, $(b \Rightarrow c)$
    
    \item Déduire $\neg a$ à partir des hypothèses : $(d \Rightarrow b)$, $(c \wedge \neg b)$, $(a \Rightarrow d)$
\end{enumerate}
\end{exercisebox}

\section{Déductions à partir de formes normales conjonctives}

\begin{exercisebox}{Exercice 11}
Pour chacune des formules, réduisez-la à la forme réduite puis utilisez les règles d'inférence afin de détailler l'argumentation permettant d'obtenir la conclusion donnée.

\begin{enumerate}
    \item $(\neg p\vee r\vee q)\wedge(b\vee c\vee d\vee\neg f)\wedge(b\vee e\vee f)\wedge(b\vee\neg b)\wedge(q\vee r\vee\neg q)\wedge(\neg p\vee q)\wedge(y\vee r\vee q)\wedge(\neg p\vee b\vee r\vee q\vee a)$
    
    Déduire $e \vee c$ à partir de cette formule.
    
    \item $(\neg b \vee a \vee r \vee b) \wedge (a \vee b \vee r \vee c) \wedge (\neg b \vee b) \wedge (\neg b \vee e \vee r \vee b) \wedge (\neg c \vee e) \wedge (e \vee b) \wedge (\neg p \vee e \vee r \vee a) \wedge (\neg p \vee e \vee r \vee b \vee a)$
    
    Déduire $b \vee e$ à partir de cette formule.
    
    \item $(a \vee c) \wedge (\neg a \vee c) \wedge (a \vee b \vee x) \wedge (b \vee \neg a) \wedge (\neg c \vee d) \wedge (a \vee b \vee \neg a) \wedge (b \vee \neg c \vee d) \wedge (\neg a \vee e \vee a)$
    
    Déduire $b \vee d$ à partir de cette formule.
    
    \item $(a \vee b \vee \neg c \vee d) \wedge (a \vee \neg b \vee \neg c) \wedge (\neg c \vee d) \wedge (\neg b \vee c) \wedge (\neg c \vee d) \wedge (a \vee b \vee c \vee d) \wedge (d \vee \neg d)$
    
    Montrez que cette formule est insatisfaisable.
    
    \item $(\neg a\vee b\vee c)\wedge(a\vee\neg a)\wedge(a\vee c\vee\neg b)\wedge(\neg b\vee\neg a)\wedge(a\vee b\vee a)\wedge(\neg b\vee b)\wedge(c\vee\neg b)\wedge(a\vee\neg c\vee\neg a)$
    
    Montrez que cette formule est insatisfaisable.
\end{enumerate}
\end{exercisebox}

\section{Modéliser et déduire}

\begin{exercisebox}{Exercice 12}
Pour chacun des cas, utilisez un formalisme logique pour traduire les hypothèses données, puis utilisez les règles d'inférence pour produire une argumentation afin de conclure.

\begin{enumerate}
    \item Il pleut ou il fait beau en ce moment, je prends ma voiture. S'il fait beau je prends mon vélo et s'il neige je prends ma voiture.
    
    Montrez qu'on peut en déduire : Je prends ma voiture, ou je prends mon vélo ou je ne sors pas.
    
    \item Si tu m'envoies tes mails, je finirai d'écrire le code. Si tu ne m'envoies pas tes mails, je me réveillerai en forme.
    
    Montrez qu'on peut en déduire : Si je ne finis pas d'écrire le code, alors je me coucher tôt. Si je ne peux pas me coucher tôt, alors je ne me réveillerai pas en forme.
    
    \item Si je fais de l'athlétisme, j'ai des courbatures. Je prends un bain quand j'ai des courbatures. Je n'ai pas pris de bain.
    
    Montrez qu'on peut en déduire : Je n'ai pas fait d'athlétisme.
    
    \item Lorsque j'hallucine, je vois des éléphants roses. Je suis en train de rêver ou j'hallucine.
    
    Montrez qu'on peut en déduire : Je bois un grand verre d'eau quand je vois des éléphants roses. Je ne rêve pas.
    
    \item Dans une maison hantée, les esprits se manifestent sous deux formes différentes : des chants discordants ou des rires moqueurs.
    \begin{itemize}
        \item Les chants ne se font pas entendre, à moins que nous jouions de l'orgue sans que les rires ne se fassent entendre.
        \item Si nous brûlons de l'encens, les rires se font entendre si et seulement si les chants se font entendre.
    \end{itemize}
    Actuellement, les chants se font entendre et les rires sont silencieux.
    
    Montrez qu'on peut en déduire : Nous jouons de l'orgue et ne brûlons pas d'encens.
\end{enumerate}
\end{exercisebox}

\section{Analyser des raisonnements}

\begin{exercisebox}{Exercice 13}
Pour chacun des cas, précisez si le raisonnement est correct en justifiant votre réponse.

\begin{enumerate}
    \item Si je réussis mes examens ou que j'ai assez d'argent, je m'offrirai une guitare. Je me suis offert une guitare.
    
    Peut-on en déduire que j'ai réussi mes examens ? Peut-on en déduire que j'ai assez d'argent ?
    
    \item Si je réussis mes examens et que j'ai assez d'argent, je m'offrirai une guitare. J'ai assez d'argent. Je ne me suis pas offert de guitare.
    
    Peut-on en déduire que je n'ai pas réussi mes examens ?
    
    \item Si je réussis mes examens ou que j'ai assez d'argent, je m'offrirai une guitare. J'ai réussi mes examens.
    
    Peut-on en déduire que je me suis offert de guitare ?
    
    \item Si je réussis mes examens ou que j'ai assez d'argent, je m'offrirai une guitare. Je n'ai pas réussi mes examens et je n'ai pas assez d'argent.
    
    Peut-on en déduire que je ne me suis pas offert de guitare ?
    
    \item Si je réussis mes examens et que j'ai assez d'argent, je m'offrirai une guitare. J'ai réussi mes examens.
    
    Peut-on en déduire que si j'ai assez d'argent alors je m'offrirai une guitare ?
\end{enumerate}
\end{exercisebox}

\section{Pour aller plus loin}

\begin{exercisebox}{Exercice 14 - Le problème des Quatre Reines}
Le problème des Huit Reines est le suivant : est-il possible de placer 8 reines sur un échiquier $8 \times 8$ sans qu'aucune reine n'en menace une autre ?

Considérons une version plus simple où on souhaite placer seulement 4 reines sur un échiquier $4 \times 4$. On utilisera les variables logiques suivantes :
\begin{itemize}
    \item $x_{i,j}$ : Une reine se trouve dans la case $i, j$
\end{itemize}

On notera les cases de la façon suivante :
\begin{center}
\begin{tabular}{|c|c|c|c|}
\hline
1,1 & 1,2 & 1,3 & 1,4 \\
\hline
2,1 & 2,2 & 2,3 & 2,4 \\
\hline
3,1 & 3,2 & 3,3 & 3,4 \\
\hline
4,1 & 4,2 & 4,3 & 4,4 \\
\hline
\end{tabular}
\end{center}

\begin{enumerate}
    \item Formalisez le fait qu'exactement une reine doit se trouver sur la première ligne.
    \item Généralisez pour toutes les lignes, toutes les colonnes, toutes les diagonales.
    \item En utilisant les hypothèses établies dans les questions précédentes, montrez que si une reine est placée dans la case (1,1), alors il n'y a pas de solution.
\end{enumerate}
\end{exercisebox}

\newpage

% ============================================================================
% TD3 - LOGIQUE DE PREMIER ORDRE
% ============================================================================

\part{TD3 - Logique de premier ordre}

\section*{Objectif}
L'objectif de ce TD est de mettre en pratique les bases de la logique de premier ordre.

\section{De la logique de premier ordre au français}

\begin{exercisebox}{Exercice 1}
Traduisez en français les formules logiques suivantes où les prédicats signifient :
\begin{itemize}
    \item $E(x)$ : ``$x$ est étudiant.''
    \item $P(x)$ : ``$x$ a visité Paris.''
    \item $C(x)$ : ``$x$ a un chat.''
    \item $F(x)$ : ``$x$ a une forêt.''
\end{itemize}

\begin{enumerate}
    \item $\forall x(P(x) \wedge C(x))$
    \item $\exists x(C(x) \wedge \neg F(x))$
    \item $\forall x(E(x) \wedge P(x) \Rightarrow C(x))$
    \item $\forall x(P(x) \Rightarrow F(x))$
    \item $\neg\exists x(E(x) \wedge C(x) \wedge F(x))$
\end{enumerate}
\end{exercisebox}

\section{Du français à la logique de premier ordre}

\subsection{Traduire des phrases}

\begin{exercisebox}{Exercice 2}
Donnez la formule logique correspondant aux phrases suivantes en utilisant le prédicat suivant : 

$D(x, y)$ : ``La personne $x$ peut duper la personne $y$''

On considérera que le domaine des variables est tout le monde.

\begin{enumerate}
    \item Tout le monde peut duper Alban.
    \item Brigitte peut duper tout le monde.
    \item Tout le monde peut duper quelqu'un.
    \item Il n'y a personne qui puisse duper tout le monde.
    \item Tout le monde peut être dupé par quelqu'un.
    \item Personne ne peut duper à la fois Alban et Brigitte.
    \item Personne ne peut se duper lui-même.
    \item Il y a des gens qui peuvent duper aucune personne à part eux-mêmes.
\end{enumerate}
\end{exercisebox}

\subsection{Quelques négations}

\begin{exercisebox}{Exercice 3}
Pour chacune des phrases précédentes, donnez la négation de la formule obtenue, puis donnez la phrase en français correspondant à cette négation.
\end{exercisebox}

\section{Identifier des inférences}

\begin{exercisebox}{Exercice 4}
Pour chacun des cas, utilisez un formalisme logique pour traduire les hypothèses et la conclusion et précisez la règle d'inférence utilisée pour obtenir cette conclusion à partir des hypothèses.

\begin{enumerate}
    \item ``Tous les hommes sont mortels. Or Socrate est un homme.'' On en déduit : ``Socrate est mortel.''
    
    \item ``Alice aime les mathématiques.'' On en déduit : ``Il y a des gens qui aiment les mathématiques.''
    
    \item ``Tous les gens qui n'aiment pas les burgers aiment les haricots. Tous les gens qui aiment les haricots aiment le sport.'' On en déduit : ``Tous les gens qui n'aiment pas les burgers aiment le sport.''
    
    \item ``Tout le monde aime les clafoutis.'' On en déduit : ``Marc aime le clafoutis.''
    
    \item ``Tous les gens qui ont un ordinateur savent envoyer un mail. Mon grand-père n'a pas d'ordinateur.'' On en déduit : ``Mon grand-père ne sait pas envoyer de mails.''
\end{enumerate}
\end{exercisebox}

\section{Modéliser et déduire}

\begin{exercisebox}{Exercice 5}
Pour chacun des cas, utilisez un formalisme logique pour traduire les hypothèses données et la conclusion, puis utilisez les règles d'inférence pour procéder à une argumentation permettant de conclure.

\begin{enumerate}
    \item ``Tous les colibris ont des couleurs vives. Aucun gros oiseau ne se nourrit de miel et ont des couleurs ternes.'' Montrez qu'on peut en déduire : ``Les colibris sont petits.''
    
    \item ``Tous les gens qui font du sport mangent sainement et boivent beaucoup d'eau. Tous les étudiants de l'université font du sport et dorment bien la nuit et boivent beaucoup d'eau.'' Montrez qu'on peut en déduire : ``Tous les étudiants de l'université dorment bien la nuit et boivent beaucoup d'eau.''
    
    \item ``Tout le monde a un chat ou un oiseau. Tous les gens qui ont un oiseau mais pas de chat ont un poisson rouge.'' Montrez qu'on peut en déduire : ``Tous les gens qui n'ont pas de poisson rouge ont un chat.''
    
    \item ``Tout le monde aime voyager ou rester chez soi et écouter des musiques. Tout le monde aime faire du sport ou ne pas écouter de musique. Tous les gens qui aiment les musiques mais n'aiment pas le sport. Jeanne aime une personne qui n'aime pas voyager.'' Montrez qu'on peut en déduire : ``Il y a des gens qui n'aiment pas le chant.''
    
    \item ``Au zoo, Koko est un singe. Tous les animaux du zoo aiment jouer ou aiment les légumes. Tous les animaux qui aiment les pommes aiment les légumes. Les singes détestent le poulet.'' Montrez qu'on peut en déduire : ``Il y a des animaux qui aiment jouer dans ce zoo.''
\end{enumerate}
\end{exercisebox}

\section{Analyser des raisonnements}

\begin{exercisebox}{Exercice 6}
Pour chacun des cas, précisez si le raisonnement est correct en justifiant votre réponse.

\begin{enumerate}
    \item ``Tous les étudiants de première année suivent un cours de programmation. Alice est étudiante de première année.'' Peut-on en déduire qu'Alice suit des cours de programmation ?
    
    \item ``Tous les étudiants d'informatique suivent un cours de programmation. Bob suit des cours de programmation.'' Peut-on en déduire que Bob est étudiant d'informatique ?
    
    \item ``Tous les perroquets aiment les fruits. Tui n'est pas un perroquet.'' Peut-on en déduire que Tui n'aime pas les fruits ?
    
    \item ``Toutes celles et ceux qui mangent une pomme chaque jour sont en bonne santé. Carole n'est pas en bonne santé.'' Peut-on en déduire que Carole ne mange pas de pomme chaque jour ?
    
    \item ``Des amis aiment tous les films d'action. Dan aime le film 'Die Hard'.'' Peut-on en déduire que 'Die Hard' est un film d'action ?
\end{enumerate}
\end{exercisebox}

\section{Montrer des règles d'inférences}

\begin{exercisebox}{Exercice 7}
Nous supposons ici que l'on connaît et peut appliquer les règles d'inférence de la logique propositionnelle ainsi que les règles d'instanciation universelle, de généralisation universelle, d'instanciation existentielle et de généralisation existentielle.

\begin{enumerate}
    \item Montrez que la règle de modus ponens universel est correcte.
    \item Montrez que la règle de modus tollens universel est correcte.
    \item Montrez que la règle de transitivité universelle est correcte.
\end{enumerate}
\end{exercisebox}

\section{Algorithmiques et quantificateurs}

\begin{exercisebox}{Exercice 8}
On donne le pseudo-code de deux procédures p1 et p2 :

\begin{minipage}{0.45\textwidth}
\begin{lstlisting}
Procedure p1()
    bool <- vrai
    PourTout i dans A Faire
        PourTout j dans B Faire
            bool <- bool ET F(i,j)
        FinPourTout
    FinPourTout
    Retourner bool
FinProcedure
\end{lstlisting}
\end{minipage}
\hfill
\begin{minipage}{0.45\textwidth}
\begin{lstlisting}
Procedure p2()
    bool <- vrai
    PourTout i dans A Faire
        bool2 <- faux
        PourTout j dans B Faire
            bool2 <- bool2 OU P(i,j)
        FinPourTout
        bool <- bool ET bool2
    FinPourTout
    Retourner bool
FinProcedure
\end{lstlisting}
\end{minipage}

\begin{enumerate}
    \item Observez le pseudo-code de la procédure p1. Donnez la formule logique qui est vraie si et seulement si p1 renvoie vrai.
    \item Observez le pseudo-code de la procédure p2. Donnez la formule logique qui est vraie si et seulement si p2 renvoie vrai.
    \item Donnez le pseudo-code de la procédure p3, qui renvoie vrai si et seulement si la formule $\exists x\forall yP(x, y)$ est vraie.
    \item Donnez le pseudo-code de la procédure p4, qui renvoie vrai si et seulement si la formule $\exists x\exists yP(x, y)$ est vraie.
\end{enumerate}
\end{exercisebox}

\newpage

% ============================================================================
% TD4 - PREUVES DE CORRECTION
% ============================================================================

\part{TD4 - Preuves de correction}

\section*{Objectif}
L'objectif de ce TD est de pratiquer les preuves de correction d'algorithmes.

\section{Procédure simple}

\begin{exercisebox}{Exercice 1}
Soit l'algorithme simple suivant :
\begin{lstlisting}
y <- 2
z <- x + y
x <- y - z
\end{lstlisting}

\begin{enumerate}
    \item Montrez la correction partielle de l'extrait ci-dessus pour l'assertion initiale $x = 5$ et l'assertion finale $x = 14$
    \item Montrez la correction partielle de l'extrait ci-dessus pour l'assertion initiale $x = -5$ et l'assertion finale $x = -6$
    \item Montrez la correction partielle de l'extrait ci-dessus pour l'assertion initiale $x = -2$ et l'assertion finale $x = 0$
\end{enumerate}
\end{exercisebox}

\section{Procédures avec blocs conditionnels}

\begin{exercisebox}{Exercice 2}
\textbf{Question 1 :} Montrez que l'extrait de procédure suivant est partiellement correct pour l'assertion initiale T et l'assertion finale $z > 0$

\begin{lstlisting}
Si x < 0 Alors
    z <- 0
FinSi
\end{lstlisting}

\textbf{Question 2 :} Montrez que l'extrait de procédure suivant est partiellement correct pour l'assertion initiale T et l'assertion finale $(x \leq y \wedge min = x) \vee (x > y \wedge min = y)$

\begin{lstlisting}
Si x < y Alors
    min <- x
Sinon
    min <- y
FinSi
\end{lstlisting}
\end{exercisebox}

\section{Procédures avec boucles}

\subsection{Calcul de puissance}

\begin{exercisebox}{Exercice 3.1}
On souhaite montrer que la procédure suivante est correcte pour l'assertion initiale ``$x \in \mathbb{R}, n \in \mathbb{N}^{**}$'' et l'assertion finale ``puissance$(x, n) = x^n$''

\begin{lstlisting}
Procedure puissance(x : reel, n : entier strictement positif)
    puissance <- 1
    i <- 0
    TantQue i < n Faire
        puissance <- puissance * x
        i <- i + 1
    FinTantQue
    Renvoyer puissance
FinProcedure
\end{lstlisting}

\begin{enumerate}
    \item Donnez un invariant de boucle
    \item Vérifiez la correction partielle de la procédure en utilisant l'invariant de boucle de la question précédente
    \item Montrez que la procédure termine et déduisez-en qu'elle est correcte
\end{enumerate}
\end{exercisebox}

\subsection{Division entière}

\begin{exercisebox}{Exercice 3.2}
On souhaite montrer que la procédure suivante est correcte pour l'assertion initiale ``$a, d \in \mathbb{N}^*$'' et l'assertion finale ``$q, r \in \mathbb{N}^*$ où $a = dq + r$ et $0 \leq r < d$''

\begin{lstlisting}
Procedure division(a, d : entiers strictement positifs)
    q <- 0
    TantQue r >= d Faire
        r <- r - d
        q <- q + 1
    FinTantQue
    Renvoyer q, r
FinProcedure
\end{lstlisting}

\begin{enumerate}
    \item Donnez un invariant de boucle
    \item Vérifiez la correction partielle de la procédure en utilisant l'invariant de boucle de la question précédente
    \item Montrez que la procédure termine et déduisez-en qu'elle est correcte
\end{enumerate}
\end{exercisebox}

\section{Procédures récursives}

\begin{exercisebox}{Exercice 4}
\textbf{Question 1 :} Montrez que la procédure \texttt{somme}$(n)$ donnée calcule $\sum_{k=0}^{n} k$

\begin{lstlisting}
Procedure somme(n : entier positif)
    Si n = 0 Alors
        Renvoyer 0
    Sinon
        Renvoyer n + somme(n - 1)
    FinSi
FinProcedure
\end{lstlisting}

\textbf{Question 2 :} Montrez que la procédure \texttt{factorielle}$(n)$ donnée calcule $n!$

\begin{lstlisting}
Procedure factorielle(n : entier positif)
    Si n = 0 Alors
        Renvoyer 1
    Sinon
        Renvoyer n * factorielle(n - 1)
    FinSi
FinProcedure
\end{lstlisting}
\end{exercisebox}

\section{Analyser des procédures}

\subsection{Procédure $p_{5,1}$}

\begin{exercisebox}{Exercice 5.1}
\begin{lstlisting}
Procedure p_{5,1}(n : entier strictement positif, m : entier)
    Si n = 1 Alors
        Renvoyer m
    Sinon
        Renvoyer m + p_{5,1}(n - 1, m)
    FinSi
FinProcedure
\end{lstlisting}

\begin{enumerate}
    \item Quel est le résultat renvoyé par la procédure $p_{5,1}$ en donnant en entrée $n = 2$ et $m = 2$
    \item Quel est le résultat renvoyé par la procédure $p_{5,1}$ en donnant en entrée $n = 3$ et $m = 0$
    \item Quel est le résultat renvoyé par la procédure $p_{5,1}$ en donnant en entrée $n = 3$ et $m = 4$
    \item Déduisez des questions précédentes ce que calcule la procédure $p_{5,1}$
    \item Montrez que la procédure $p_{5,1}$ est correcte
\end{enumerate}
\end{exercisebox}

\subsection{Procédure $p_{5,2}$}

\begin{exercisebox}{Exercice 5.2}
\begin{lstlisting}
Procedure p_{5,2}(x : reel)
    y <- 2
    z <- x + y
    Si z > 0 Alors
        z <- 3
    Sinon
        z <- 0
    FinSi
    Renvoyer z
FinProcedure
\end{lstlisting}

\begin{enumerate}
    \item Quel est le résultat renvoyé par la procédure $p_{5,2}$ en donnant en entrée $x = 3$
    \item Quel est le résultat renvoyé par la procédure $p_{5,2}$ en donnant en entrée $x = -2$
    \item Quel est le résultat renvoyé par la procédure $p_{5,2}$ en donnant en entrée $x = -4$
    \item Déduisez des questions précédentes ce que calcule la procédure $p_{5,2}$
    \item Montrez que la procédure $p_{5,2}$ est correcte
\end{enumerate}
\end{exercisebox}

\section{Identifier des règles d'inférences}

\subsection{Pour blocs conditionnels}

\begin{exercisebox}{Exercice 6.1}
En utilisant la seconde règle d'inférence pour la correction partielle des blocs conditionnels vue en cours, proposez une règle d'inférence pour prouver la correction partielle du bloc conditionnel suivant :

\begin{lstlisting}
Si c1 Alors
    S1
Sinon Si c2 Alors
    S2
Sinon
    S3
FinSi
\end{lstlisting}

\textbf{Question 2 :} Généralisez pour un bloc conditionnel de la forme :

\begin{lstlisting}
Si c1 Alors
    S1
Sinon Si c2 Alors
    S2
Sinon Si c3 Alors
    S3
Sinon
    Sn
FinSi
\end{lstlisting}
\end{exercisebox}

\subsection{Pour des cas quelconques}

\textbf{Question 1 :} On suppose que les hypothèses suivantes sont vraies :
\begin{itemize}
    \item $q_0 \Rightarrow q_1$
    \item $p\{S\}q_1$
\end{itemize}
Montrez qu'on peut en déduire $p\{S\}q$

\textbf{Question 2 :} On suppose que les hypothèses suivantes sont vraies :
\begin{itemize}
    \item $p_0 \Rightarrow p_1$
    \item $p_1\{S\}q$
\end{itemize}
Montrez qu'on peut en déduire $p_0\{S\}q$

\vspace{2cm}

\begin{center}
\rule{10cm}{0.5pt}

\vspace{0.5cm}

\textit{Fin des Travaux Dirigés}

\vspace{0.5cm}

\textbf{CERI - Avignon Université}

\textbf{Mathématiques Discrètes - Licence 1 Informatique}

\textbf{Hiver 2026}
\end{center}

\end{document}